\section{Why the elementary proofs fit iPRA}
\label{sec:arithmetic-audit}

This section explains the claimed weak-arithmetic status of the written
proofs. It does not identify native Lean induction with induction inside
iPRA. The intended object theory is the conventional first-order theory
implemented by ProveIt's \texttt{PRAAx}.

\subsection{Language, axioms, and rules}

The underlying domain is $\N$. There is one function symbol for each
primitive recursive construction: zero, successor, projections,
composition, and primitive recursion. For $f$ of arity $k$ and $g$ of
arity $k+2$, the recursion symbol obeys
\[
 \begin{split}
 h(0,\vec u)&=f(\vec u),\\
 h(n+1,\vec u)&=g(n,h(n,\vec u),\vec u).
 \end{split}
\]
Composition and projections have their usual defining equations.
The arithmetic axioms include $Sx\ne0$ and $Sx=Sy\to x=y$.
For each quantifier-free formula $\theta(n,\vec u)$, the induction
schema is
\[
 \theta(0,\vec u)\to
 \bigl[\forall n\,(\theta(n,\vec u)\to\theta(n+1,\vec u))\bigr]
 \to\forall n\,\theta(n,\vec u).
\]
The logic is intuitionistic first-order natural deduction with equality:
the introduction and elimination rules for $\wedge,\vee,\to,\forall,\exists$,
falsity elimination, equality reflexivity, and substitution of equals.
There is no unrestricted excluded middle or double-negation elimination.
Decidable arithmetic cases, such as the sign of an integer, are derived
using primitive recursive characteristic functions.

Signed integers use the canonical natural-number encoding in
Section~\ref{sec:logical-audit}. The integer ring laws, order laws,
absolute-value identities, and powers translate into equations and
inequalities between these codes. Their familiar elementary proofs use
only quantifier-free induction. Writing $x\in\Z$ in the arithmetic proofs
abbreviates this representation, not a new unbounded sort or quotient
principle added to the object theory.

\subsection{Concrete finite data}

Use dense coefficient arrays in increasing degree order. Normalize an
array by deleting trailing zero coefficients, with the empty array as
the zero polynomial. Length, entry lookup with default zero, truncation,
normalization, coefficient sums, addition, scalar multiplication, and
convolution are primitive recursive operations on finite sequence codes.
Horner evaluation is the recursion
\[
 E([],x)=0,\qquad E(a::L,x)=a+xE(L,x).
\]
Every polynomial identity in a certificate is checked over a specified
finite coefficient range. With array lengths $l_f,l_A,l_R$, it suffices
to test indices below $\max(l_f,2l_A,l_R)$. Outside that range both
sides are zero by their support bounds. This is why the certificate
predicate is primitive recursive rather than a universally quantified
semantic identity.

Rational arithmetic can be represented by a signed numerator and a
strictly positive denominator. Reduction to lowest terms is optional.
Addition and multiplication use the usual cross-multiplication formulas;
equality is checked by cross multiplication. If canonical rationals are
desired, the greatest common divisor can be computed by scanning divisors
up to the smaller nonzero input, with explicit finite fuel. None of the
proofs needs an unbounded search for a rational square root.

\subsection{Induction for coefficient estimates}

Fix a degree bound $m\ge1$, an integer $x$, and $t=|x|\ge1$.
For $k\le m$, define
\[
 T_k=\sum_{i<k}a_ix^i,\qquad S_k=\sum_{i<k}|a_i|.
\]
The required lower-coefficient estimate follows from the invariant
\[
 k\le m\quad\longrightarrow\quad |T_k|\le S_k t^{m-1}.
\]
The base case is $T_0=S_0=0$. At stage $k<m$, the triangle inequality
and $t^k\le t^{m-1}$ give
\[
 |T_{k+1}|\le |T_k|+|a_k|t^k
             \le(S_k+|a_k|)t^{m-1}.
\]
For $k>m$ the implication is inactive; alternatively, one freezes
the state after $m$ steps. All sums, powers, comparisons, and
absolute values are terms of primitive recursive functions. Thus the
displayed invariant is quantifier-free after definitions are expanded.
Monotonicity of powers for base $t\ge1$ follows by the same kind of
equational induction; no induction over an arbitrary family of formulas
is used.

The corresponding estimate for $R$ uses its computed leading index and
the same invariant. Its root bound is a finite sign/order argument at a
given $x$, not a search for all roots. The factor lemma uses only
\[
 |uv|=|u||v|,\qquad u\ne0\Rightarrow |u|\ge1,
 \qquad |s+t|\le|s|+|t|.
\]
When arguing by contradiction that $|x|\le B_x$, we use decidability
of that integer inequality: either it holds, or the explicit alternative
$|x|>B_x$ gives a contradiction. This is intuitionistically valid.

\subsection{Induction for square completion}

Let $u>0$ be the supplied integer square root of the leading coefficient
of $f$. The constructor has exactly $m$ descending coefficient stages
after the leading coefficient of $Q$ is set to $u$. Its state contains
a finite rational array $Q_k$, frozen after stage $m$. For $0\le k\le m$,
after $k$ stages the leading $k+1$ coefficients of $Q_k^2$ agree with
those of $f$. Encode this as a guarded bounded coefficient test
\[
 k\le m\quad\longrightarrow\quad
 \prod_{j=0}^{k}
  \operatorname{eq}_{\Q}\bigl([X^{2m-j}]Q_k^2,
                              [X^{2m-j}]f\bigr)=1.
\]
The product is an iterated Boolean conjunction. Its defining recursion
makes the assertion an atomic equality of primitive recursive terms.
The new coefficient occurs linearly with nonzero multiplier $2u$;
cross multiplication verifies the update. A second invariant states
that coefficients fixed in earlier stages are unchanged. Each step
therefore has a quantifier-free proof.

Representing denominators by products, rather than least common
multiples, avoids any issue about finding a minimum. The resulting
denominator $D$ and integer polynomial $A=DQ$ are primitive recursive
in the input array and the supplied positive square root $u$ of its
leading coefficient. Multiplying through the finite rational identities
proves $D^2f=A^2+R$. The test that the residual array is nonzero is
finite. If it is zero, the routine reports the square case and makes
no finiteness claim. Testing whether the leading coefficient is a
positive square itself requires only scanning up to that coefficient.

Uniform primitive recursiveness here refers to this \emph{fixed}
algorithm on variable input data. It is not a claim that one can
evaluate every coded primitive recursive program by a single primitive
recursive universal evaluator. There is no such general evaluator.
The implementation task is to construct one finite \texttt{PR} expression
for each fixed operation, then compose them.

\subsection{Exhaustive enumeration and its invariant}

Once $B_x$ is computed, one may scan $x=-B_x,\ldots,B_x$. For each
$n=f(x)\ge0$, compute the greatest $s\le n$ with $s^2\le n$ by a
bounded scan, or use a separately justified fast integer-square-root
algorithm. Retain $(x,s)$ and, for $s>0$, $(x,-s)$ precisely when
$s^2=n$. Every step has an explicit bound. This implementation avoids
the larger rectangular scan using $B_y$, while the latter remains a
simple proof-theoretic specification.

After $k$ horizontal positions have been processed, the retained list
contains exactly the solutions at those positions. In the rectangular
version, encode this as the conjunction over the first $k$ box indices
that each successful point test has its decoded tuple in the accumulated
list, and that every accumulated tuple passes the test and has a
processed index. All these quantifiers are bounded by the current code
lengths or the explicit box size. Their characteristic functions turn
the invariant into a quantifier-free formula. The single update either
appends the next successful tuple or leaves the list unchanged, which
proves the next invariant by finite membership identities.

At termination, the invariant proves exactness within the box. The
Runge theorem independently proves that every integral point lies in
the box. Combining these two conclusions proves global completeness.
The program is therefore an exhaustive solver on the stated certified
family, not just an experimental search below a chosen cutoff.

\subsection{What remains to be formalized}

The written arguments give the elementary arithmetic proof and its
allowed induction invariants. Completing the embedded derivation still
requires coding those functions as actual \texttt{PR} terms, deriving
their equations in \texttt{PRAAx}, constructing the characteristic
formulas and their \texttt{QF} witnesses, and assembling natural-deduction
proof terms. These are concrete finite tasks for the elementary family.
Their completion is not asserted by the presence of a native Lean
specification or by the Python checks.

The external identification of the examples' projective genus uses
standard algebraic geometry. The integral-point theorem itself does
not invoke that identification as an axiom: its hypotheses and
conclusion are directly arithmetic. A separate formalization of the
geometric interpretation would be needed to make the genus statement
part of the same internal arithmetic development.
